\chapter{Herpes Simplex}
\index{Herpes Simplex}

\section*{Definition and Overview}
Herpes simplex is a viral infection caused by herpes simplex virus type 1 (HSV-1) and type 2 (HSV-2). HSV-1 typically causes cold sores around the mouth and face, while HSV-2 typically causes genital herpes, though both types can infect either site. After the first infection, the virus remains dormant in nerve cells and can reactivate, causing recurrent sores. Infection is extremely common, often produces no symptoms, and is transmitted by close contact. Antiviral medicines manage outbreaks effectively, and serious complications are rare in healthy people.

\section*{Epidemiology}
Herpes simplex is one of the most prevalent infections in the world. Around 3.7 billion people under 50 have HSV-1, and around 500 million have HSV-2. Most infected people are unaware because symptoms are absent or mild. In many countries, around 70--80\% of adults have antibodies to HSV-1, and approximately one in eight has HSV-2. First infections often occur in childhood through non-sexual contact. Recurrences are common but tend to decrease with time.

\section*{Aetiology and Pathophysiology}
\begin{itemize}
    \item \textbf{HSV-1:} transmitted by kissing, sharing utensils or lip balm, and oral sex; causes oral cold sores and some genital infections
    \item \textbf{HSV-2:} transmitted sexually; the main cause of genital herpes
    \item \textbf{Transmission:} contact with skin shedding virus, including when no sores are visible
    \item \textbf{Primary infection:} the virus enters through skin or mucous membranes and replicates locally, causing blisters
    \item \textbf{Latency:} the virus retreats to the sensory ganglia near the spine and remains dormant for life
    \item \textbf{Reactivation:} triggers including stress, illness, fever, sun exposure, hormonal changes, and immunosuppression cause recurrent outbreaks
    \item \textbf{Shedding:} asymptomatic shedding means the virus can be passed on without visible sores
\end{itemize}

\section*{Clinical Features}
\begin{itemize}
    \item Cold sores: tingling, then clusters of small blisters on or around the lips, crusting and healing within 7--10 days
    \item Genital herpes: painful blisters and sores on the genitals, buttocks, or thighs
    \item Primary infection: may include fever, swollen glands, and flu-like symptoms
    \item Gingivostomatitis: widespread mouth ulcers with fever, mainly in young children with first infection
    \item Herpetic whitlow: painful infection of a finger, often in health workers
    \item Recurrences are shorter, milder, and often preceded by tingling
    \item Eye infection (herpes keratitis) causes pain, redness, and sensitivity to light
\end{itemize}

\section*{Differential Diagnosis}
\begin{itemize}
    \item Aphthous ulcers and other mouth conditions for oral sores
    \item Syphilis, chancroid, and other causes of genital ulcers
    \item Herpes zoster (shingles), which has a similar blistering appearance
    \item Impetigo and other skin infections
    \item Contact dermatitis and eczema herpeticum
    \item PCR swabs confirm herpes simplex when the diagnosis is unclear
\end{itemize}

\section*{When to Seek Medical Care}
See a doctor for suspected genital herpes, frequent or severe cold sores, or sores that do not heal. Seek urgent care for eye symptoms, severe infection, or symptoms in someone with a weakened immune system.

\textbf{Contact your local emergency services for:}
\begin{itemize}
    \item Severe headache, confusion, or seizures with herpes (possible encephalitis)
    \item Eye pain, redness, or blurred vision with a cold sore (possible keratitis)
    \item A newborn with fever, rash, or lethargy whose mother has herpes
    \item Difficulty breathing or signs of a severe reaction
    \item Extensive or severe infection in an immunocompromised person
\end{itemize}

\section*{Diagnosis and Investigations}
\begin{itemize}
    \item \textbf{Clinical diagnosis:} typical blisters and sores are usually recognisable
    \item \textbf{PCR swab:} the most sensitive test, confirming the virus and type
    \item \textbf{Viral culture:} older method, less sensitive
    \item \textbf{Blood antibody tests:} detect past infection and type
    \item \textbf{Eye assessment:} slit-lamp examination for suspected eye involvement
    \item \textbf{CSF analysis:} for suspected encephalitis
\end{itemize}

\section*{Management}
\begin{itemize}
    \item \textbf{Topical antivirals:} creams such as aciclovir, most useful when started early
    \item \textbf{Oral antivirals:} aciclovir, valaciclovir, or famciclovir for first episodes and severe recurrences
    \item \textbf{Suppressive therapy:} daily oral antivirals for frequent recurrences
    \item \textbf{Symptom relief:} analgesics, cool compresses, and lip balms with sun protection
    \item \textbf{Good hygiene:} handwashing and avoiding touching sores to prevent spread to the eyes and others
    \item \textbf{Eye infection:} urgent ophthalmology treatment with antiviral drops and tablets
    \item \textbf{Pregnancy:} specialist care to prevent neonatal infection
\end{itemize}

\section*{Complications}
\begin{itemize}
    \item Herpes keratitis, which can scar the cornea and threaten sight
    \item Encephalitis, a rare but serious brain infection
    \item Neonatal herpes, transmitted at birth, with severe consequences
    \item Eczema herpeticum, a widespread severe rash in people with eczema
    \item Whitlow and secondary bacterial infection of sores
    \item Increased HIV risk with genital sores
    \item Psychological distress and stigma
\end{itemize}

\section*{Prognosis and Outcomes}
Herpes simplex infection is lifelong, but its impact usually diminishes over time. Recurrences become less frequent and milder, and suppressive therapy controls symptoms effectively for those with frequent outbreaks. Serious complications are rare in healthy people, and pregnancy outcomes are excellent with specialist management. Cold sores are a nuisance for most people but heal without scarring. Research into vaccines and curative therapies is progressing.

\section*{Prevention and Self-Care}
\begin{itemize}
    \item Avoid kissing and sharing utensils when cold sores are active
    \item Use condoms to reduce genital transmission
    \item Use lip balm with sunscreen to reduce sun-triggered cold sores
    \item Manage stress and illness to reduce recurrences
    \item Wash hands after touching sores
    \item See a doctor promptly for eye symptoms
    \item Pregnant women with herpes should inform their maternity team
    \item Start antiviral treatment early in an outbreak for best effect
\end{itemize}

\section*{Sources and Further Reading}
The following reputable sources provide further information for healthcare students and patients seeking reliable, current advice about herpes simplex.

\begin{itemize}
    \item Healthdirect. \textit{Herpes simplex virus: symptoms and treatment.} https://www.healthdirect.gov.au/herpes-simplex-virus
    \item Melbourne Sexual Health Centre. \textit{Herpes simplex information.} https://www.mshc.org.au/
    \item World Health Organization. \textit{Herpes simplex virus: fact sheet.} https://www.who.int/
    \item NHS. \textit{Cold sores: overview and treatment.} https://www.nhs.uk/conditions/cold-sores/
    \item Mayo Clinic. \textit{Cold sores: symptoms and causes.} https://www.mayoclinic.org/
    \item MedlinePlus. \textit{Herpes simplex.} https://medlineplus.gov/herpessimplex.html
\end{itemize}
